\section{Petri/CE/PT Nets}
In contrast to hierarchical state machines, state transitions in Petri Nets are asynchronous. Focus on modelling causal dependencies; no global synchronization assumed \f{\ \rightarrow\ } message passing only. Petri nets can be used to model concurrent distributed systems (e.g. UML).
\begin{figure}[h!]
    \centering
    \includegraphics[width=.5\textwidth]{petri.png}
\end{figure}

\subsection{Components}
\begin{itemize}
    \item \b{Conditions:} Represent "local states" and are either met or not met (potential state space)
    \item \b{Events:} May take place if certain conditions are met (represents state transition)
    \item \b{Flow relation:} Relates conditions and events
    \item \b{Tokens:} Assignment of tokens to conditions specify a global state 
\end{itemize}
\definition{Petri Net}{
    A Petri Net \f{N = (C,E,F,M_0)} is defined with:
    \begin{enumerate}
        \item \f{C}: set of conditions/places
        \item \f{E}: set of events/transitions
        \item \f{F}: set of edges (binary flow relations) with \f{F\subseteq(C\times E)\cup(E\times C),\quad}\f{C} and \f{E} disjoint
        \item \f{M_0: C\rightarrow \mathbb{N}} as the initial marking
    \end{enumerate}
} 
\definition{Preconditions / Postconditions}{
    Let \f{N} be a Petri Net and let \f{x\in E, y\in C}:
    \begin{enumerate}
        \item \f{\cdot x := \left\{y\ |\ (y\times x)\in F\right\}} is called the set of preconditions
        \item \f{x\cdot := \left\{y\ |\ (x\times y)\in F\right\}} is called the set of postconditions
    \end{enumerate}
}

\subsection{Marking}
A marking is a mapping \f{M:C\rightarrow\left\{0,1\right\}}. Firing event \f{x} generates new markings on all conditions \f{C} according to the following rules:
\begin{enumerate}
    \item An event \f{x} can be fired, iff \f{\cdot x} is marked and \f{x\cdot} is unmarked
    \item After firing, \f{\cdot x} is unmarked and \f{x\cdot} is marked
    \item \f{M\rightarrow M'}, iff \f{M'} results from \f{M} by firing \it{exactly} one transition
\end{enumerate}

\subsection{Condition-Event Nets}
\definition{Different Concepts}{
    \begin{enumerate}
        \item \f{(c,e)\in C\times E\rightarrow (c,e)} is called \b{loop} iff \f{cFe\wedge eFc}.
        \item Net \f{N=(C,E,F)} is called \b{pure}, if \f{F} does not contain loops.
        \item Net is called \b{simple}, if there are no two transitions that have the same sets of input and output places: \f{\forall x,y\in E: \cdot x = \cdot y \wedge x\cdot = y\cdot \ \Rightarrow\  x=y}
    \end{enumerate}
}
\definition{C/E-Nets}{
    A Petri Net \f{N=(C,E,F)} together with a non-empty set of markings \f{M} is called C/E-Net, iff:
    \begin{itemize}
        \item \f{N} is simple, pure, ad has no isolated elements
        \item \f{M} is closed wrt. "firing" and "inverse firing"
        \item \f{N} is deadlock-free: For each event \f{e\in E} exists a marking \f{m\in M} that allows firing at \f{e}
    \end{itemize}
}
\subsection{Properties}
\begin{itemize}
    \item Marking \f{M'} is \b{reachable} from marking \f{M}, iff there exists a sequence of firings that transform \f{M} into \f{M'}
    \item C/E net is \b{cyclic}, iff any two markings are reachable from each other
    \item C/E net fulfills \b{liveness}, iff \f{N} is deadlock-free
    \begin{itemize}
        \item \b{Strong liveness:} For each \f{M} and \f{e} exists \f{M'} so that \f{M'} is reachable from \f{M} and \f{M'} activates \f{e} for firing.
        \item \b{Weak liveness:} For each \f{M} an event \f{e} exists so that \f{M} activates \f{e} for firing.
    \end{itemize}
\end{itemize}

\subsection{Place/Transition Nets}
PT-Nets can model resource (i.e. producer/consumer) usage and causal relationships, by allowing more than one token per condition, capacities of places, and weights of edges. The downside is that the state space may become infinite.\\
Nets can be used to verifiy properties such as reachability, liveness, or bounds on the number of tokens.

\definition{PT-Nets}{
    A net \f{N=(P,T,F,K,W,M_0)} is called PT-Net (non-deterministic), iff:
    \begin{itemize}
        \item \f{N=(P,T,F)} is a net with places \f{p\in P} and transitions \f{t\in T}
        \item \f{K:P\rightarrow (N_0\cup\left\{\omega\right\})\backslash \left\{0\right\}} denotes the capacity of places; \f{\omega\equiv \infty} capacity. A place \f{p_i} can hold \f{K(p_i)} tokens at maximum.
        \item \f{W: F\rightarrow (N_0\backslash \left\{0\right\})} denotes the weight of graph edges
        \item \f{M_0: P\rightarrow N_0\cup\left\{\omega\right\}} represents the initial marking of places
    \end{itemize}
    A net \f{N} is \b{contact-free}, if capacities in \f{\cdot t} never prevent activation of the transition \f{t}. A \b{pure} and \b{contact-free} PT-Net is completely defined by \f{N} and \f{M_0}.\\
    Since capacities of places are not considered in \f{N} (here: matrix \f{P\times T}), we can not always compute transitions!
}
\begin{figure}[h!]
    \centering
    \includegraphics[width=.7\textwidth]{ptnet.png}
\end{figure}

\subsubsection{Pros and Cons}
\begin{table}[h!]
\begin{tabular}{|p{.45\textwidth}|p{.475\textwidth}|}\hline
\textbf{Andvantages}                                          & \textbf{Disadvantages}                           \\\hline\hline
Can model distributed applications & problems with modeling timing \\\hline
Well-known theory for proving properties & no programming elements \\\hline
                                 & no hierarchy \\\hline
\end{tabular}
\end{table}










